\documentclass{biometrika}
\usepackage{float}
\usepackage[ruled]{algorithm2e}
\SetKwInOut{Input}{Input}
\SetKwInOut{Output}{Output\,}
\SetKwInOut{Data}{Data}
\SetKwProg{Decomp}{Decomp}{}{EndDecomp}
\SetKwProg{CombineIndices}{PartitionIndices}{}{EndPartitionIndices}
\SetKwProg{Vstatistics}{V statistics}{}{End V statistics}
\usepackage{newtxtext}
\usepackage[subscriptcorrection]{newtxmath}
\usepackage{enumitem}
\usepackage{threeparttable}
\usepackage[page]{appendix}
\usepackage{graphicx}
\usepackage{bibunits}
\usepackage{booktabs}

\renewcommand{\appendixpagename}{\large\centering Supplementary Materials for \\ ``\mytitle''}

\usepackage{natbib}

\graphicspath{{./Figures/}}

\usepackage{subcaption}

%%% User-defined macros should be placed here, but keep them to a minimum.
\def\Bka{{\it Biometrika}}
\def\AIC{\textsc{aic}}
\def\T{{ \mathrm{\scriptscriptstyle T} }}
\def\v{{\varepsilon}}

%\addtolength\topmargin{35pt}
\DeclareMathOperator{\Thetabb}{\mathcal{C}}

\newcommand{\Lin}[1]{{\color{red}{Lin: #1}}}
\newcommand{\Chen}[1]{{\color{blue}{Chen: #1}}}
\renewcommand{\_}{\raisebox{.5pt}{\underline{\phantom{x}}}}

\def\expit{\mathrm{expit}}
\def\logit{\mathrm{logit}}
\def\CI{\mathrm{CI}}
\def\det{\mathsf{det}}
\def\fam{\mathsf{fam}}
\def\Frechet{Fr\'{e}chet}
\def\dist{\mathsf{d}}
\def\T{\mathrm{T}}
\def\Exp{\mathrm{Exp}}
\def\Log{\mathrm{Log}}
\def\Law{\mathsf{Law}}
\def\Hess{\mathsf{H}}
\def\Diff{\mathrm{D}}
\def\P{\mathrm{P}}
\def\fd{\mathsf{fd}}
\def\id{\mathsf{id}}

\def\sff{\mathsf{rd}}

\def\bias{\mathsf{bias}}

\def\bmu{\bm{\mu}}

\def\adj{\mathsf{adj}}
\def\db{\mathsf{db}}
\def\unadj{\mathsf{unadj}}
\def\dc{\mathsf{dc}}
\def\proj{\mathsf{proj}}
\def\CRE{\mathsf{CRE}}
\def\lin{\mathsf{lin}}
\def\zhang{\mathsf{zhang}}
\def\chen{\mathsf{chen}}

\def\aipw{\mathsf{aipw}}
\def\ipw{\mathsf{ipw}}
\def\plim{\mathrm{plim}}

\def\bbU{\mathbb{U}}
\def\bbV{\mathbb{V}}

\usepackage{xcolor, graphicx, appendix}
\usepackage[bookmarks, plainpages = false, colorlinks = true, citecolor = teal, linkcolor = chicago-maroon, anchorcolor = red, urlcolor = chicago-maroon]{hyperref}
\definecolor{chicago-maroon}{RGB}{128,0,0}
\definecolor{northwestern-purple}{RGB}{82,0,99}
\usepackage{url}\urlstyle{rm}
\usepackage{bm, bbm}
\usepackage{mathtools}

\newcommand{\ones}{\mathbf 1}
\newcommand{\zeros}{\mathbf 0}
\newcommand{\reals}{{\mathbb{R}}}
\newcommand{\integers}{{\mathbb{Z}}}
\newcommand{\naturals}{{\mathbb{N}}}
\newcommand{\rationals}{{\mathbb{Q}}}
\def\bSigma{\bm{\Sigma}}

\def\bOmega{\bm{\Omega}}

\def\trace{\mathsf{tr}}
\def\bbP{\mathbb{P}}
\def\bx{\mathbf{x}}
\def\v{\mathrm{v}}
\def\bbR{\mathbb{R}}
\def\bbE{\mathbb{E}}
\def\sff{\mathsf{rd}}
\def\bH{\mathbf{H}}
\def\bI{\mathbf{I}}
\def\var{\mathsf{var}}
\def\cov{\mathsf{cov}}
\def\bM{\mathbf{M}}
\def\bbeta{\bm{\beta}}
\def\op{\mathsf{op}}
\def\bA{\mathbf{A}}
\def\bv{\mathbf{v}}
\def\bX{\mathbf{X}}
\def\by{\mathbf{y}}
\def\bo{\mathbf{o}}
\def\bO{\mathbf{O}}
\def\bt{\mathbf{t}}
\def\b{\mathrm{b}}
\def\y{\mathrm{y}}
\def\IIFF{\mathbb{IF}}
\def\N{\mathrm{N}}
\def\calF{\mathcal{F}}
\def\diff{\mathrm{d}}
\def\diag{\mathsf{diag}}
\def\tilde{\widetilde}
\def\hat{\widehat}
\def\calO{\mathcal{O}}
\def\calS{\mathcal{S}}
\def\calI{\mathcal{I}}
\def\bbB{\mathbb{B}}
\def\bbX{\mathbb{X}}
\def\calX{\mathcal{X}}


\def\mytitle{Algorithm of U-Statistics 
 of Kernel Size 2 and Its Computational Complexity}

\newcommand\indep{\protect\mathpalette{\protect\independenT}{\perp}}
\def\independenT#1#2{\mathrel{\rlap{$#1#2$}\mkern2mu{#1#2}}}
\def\ci{\perp\!\!\!\perp}

\begin{document}
\title{\mytitle}
\maketitle

\begin{bibunit}[biometrika]

\section{Notations and Framework of Algorithm}
\subsection{U-statistics of kernel size 2}
We consider a class of U-Statistics whose kernel size is 2.
\begin{definition}
    Assume that $\bX_1, \ldots, \bX_n$ are samples from a
    distribution $\bbP$ on $\bbX$. $h_k: \bbX \times \bbX \to \bbR, \, k = 1,\ldots, m-1$
    is a family of kernel functions over $\bbX \times \bbX$.
    A $m-th$ order U-statistics with kernel functions $\{h_k\}_{k=1}^{m-1}$ and samples $\{\bX_k\}_{k=1}^n$ is defined by
    \begin{align*}
        \bbU_m^{(2)}(\{h_k\}_{k=1}^{m-1},\{\bX_k\}_{k=1}^n) = \frac{1}{\prod_{k=0}^{m-1} (n-k)} \sum_{\substack{\forall p \neq q, i_p \neq i_q \\ 1 \leq i_p, i_q \leq n}} h_1(\bX_{i_1}, \bX_{i_2}) h_2(\bX_{i_2}, \bX_{i_3}) \cdots h_{m-1}(\bX_{i_{m-1}}, \bX_{i_m}).
    \end{align*}
    The family of matrix $\bA^{(k)}: \bA^{(k)}_{ij} = h_k(\bX_i, \bX_j) , \, 1 \leq k \leq m-1$ are called kernel matrices of this U-statistics.
    Then the U-statistics can be rewritten as
    \begin{align*}
        \bbU_m^{(2)}(\{h_k\}_{k=1}^{m-1},\{\bX_k\}_{k=1}^n) = \frac{1}{\prod_{k=0}^{m-1} (n-k)} \sum_{\forall p \neq q, i_p \neq i_q} \bA^{(1)}_{i_1,i_2} \bA^{(2)}_{i_2,i_3} \cdots \bA^{(m-1)}_{i_{m-1}, i_m}
    \end{align*}
\end{definition}

Given the family of kernel matrices, $\bbU_m^{(2)}(\{h_k\}_{k=1}^{m-1},\{\bX_k\}_{k=1}^n)$ will be abbreviated as $\bbU_m$ for simplicity.

Our algorithm is based on the decomposition of U-statistics into a linear combination of V-statistics and the whole algorithm contains two steps:
\begin{itemize}
    \item[1.] Decomposition of U-statistics to V-statistics and corresponding coefficient.
    \item[2.] Computation of each related V-statistics.
\end{itemize}

\subsection{V-statistics with restrictions}
As U-statistics can be decomposed into a linear combination of V-statistics,
we need to establish the form of all possible V-statistics related to a U-statistics $\bbU_m$.

\begin{definition}
    Let $\mathcal{I}_m$ denote a symbol set containing $m$ different symbols and $i_1, i_2, \ldots, i_m$ is some symbols in $\calI_m$.
    $[m] = \{1, 2, \ldots, m\}$ denote the set of $m$ positive integers. Let $\calS^{(m)} = \{s_k\}_{k=1}^{|\calS^{(m)}|}$ be a partition of set $[m]$ which means $s_p \cap s_q = \emptyset$ and $\bigcup_k{s_k} = \mathcal{I}_m$.
    Then a V-statistics with restriction $\calS^{(m)}$ related to the $m$-th order U-statistics $\bbU_m$ is given by
    \begin{align*}
        \bbV_{m}^{\mathcal{S}^{(m)}} = \sum_{\substack{i_p = i_q, \forall p, q \in s_k, \\ \forall s_k \in \mathcal{S}^{(m)}}} \bA^{(1)}_{i_1,i_2} \bA^{(2)}_{i_2,i_3} \cdots \bA^{(m-1)}_{i_{m-1}, i_m}
    \end{align*}
\end{definition}

\begin{remark}
    The restriction $\calS^{(m)}$ is a partition of $[m]$ and the number of elements in each subset $s_k$ is at least 1. The restriction $\calS^{(m)}$ is used to restrict the indices of the kernel matrices in the V-statistics.
\end{remark}

\begin{remark}
    In fact, the V-statistics $\bbV_{m}^{\mathcal{S}^{(m)}}$ is uniquely determined by the partition $\calS^{(m)}$ given the kernel matrices $\{\bA^{(k)}\}_{k=1}^{m-1}$ but is independent of symbol set $\calI_m$ i.e. the $\calI_m$ can be any $m$-element any set has at least $m$ elements. $\calS^{(m)}$ is called the partition representation of a V-statistics with restriction. 
\end{remark}

\begin{remark}
    Given a symbol set $\calI_m$, a V-statistics with restriction $\calS^{(m)}$ related to the $m$-th order U-statistics $\bbU_m$ is uniquely determined by a sequence of symbol sequences over $\calI_m$ as $\mathcal{L}(\calS^{(m)}) = \{i_1i_2, i_2i_3, \ldots, i_{m-1}i_m |\, i_p = i_q, \forall p, q \in s, \forall s \in \calS^{(m)}, \, i_p, i_q \in \calI_m \}$. This called a sequence representation of $\bbV_{m}^{\mathcal{S}^{(m)}}$ over $\calI_m$. 
\end{remark}

\begin{remark}
    The sequence representation establishes a relationship between V-statistics $\bbV_{m}^{\mathcal{S}^{(m)}}$ and an undirected graph
    with loops and multiple edges, precisely, an Euler path or Euler loop.
    All different symbols in $\mathcal{L}(\calS^{(m)})$ are nodes of the graph and all subsequences of $\mathcal{L}(\calS^{(m)})$ are edges. And our computational algorithm and its complexity analysis will be based on this relationship.
\end{remark}
% For the first step, we claim our main theorem as follows

% \begin{theorem}
%     Let $U_m$ be a $m$-th order U-statistics with kernel matrices $\{\bA_k\}_{k=1}^{m-1}$. Then $U_m$ can be decomposed into a linear combination of V-statistics
% \end{theorem}

% At first we make the convention for the computational form of V-statistics related to computing a $m$-th order U-statistic with former form. 

\section{Algorithm}
\subsection{Decomposition of U-statistics to V-statistics}
Now we give our main result of the decomposition of U-statistics to V-statistics as following.

\begin{theorem}
    Let $\mathcal{I}_m$ is a $m$-element symbol set and $\{\mathcal{S}^{(m)} = \{s_k\}_{k=1}^{|\mathcal{S}^{(m)}|}\}$ are its all possible partitions. The $m$-th order U-statistics is a linear combination of upper V-statistics
    \begin{align*}
        \bbU_m = \sum_{\mathcal{S}^{(m)}} w_{\mathcal{S}^{(m)}}\bbV_m^{\mathcal{S}^{(m)}}
    \end{align*}
    where
    \begin{align*}
        w_{\mathcal{S}^{(m)}} = (-1)^{\sum_{k=1}^{|\mathcal{S}^{(m)}|} (|s_k| - 1)} \prod_{k=1}^{|\mathcal{S}^{(m)}|} (|s_k| - 1)!
    \end{align*}
\end{theorem}
\subsection{Computation of V-statistics}

\section{Computational Complexity}


We will give a graph formulization of all $V$-statistics to capture the property of complexity.

\begin{theorem}[Exact Complexity Bound]
\label{thm:bound_2}
    For a function $f(k): \naturals \to \naturals$ defined as:
    \begin{align}
        f(k) = 
        \begin{cases}
            \binom{k}{2} + \frac{k}{2} - 1, & \text{if } k \text{ is even}, \\
            \binom{k}{2},                    & \text{if } k \text{ is odd},
        \end{cases}
    \end{align}
    the computational complexity exponent of $\bbU_{m}^{\mathsf{all}}$ is characterized by:
    \begin{align}
        \operatorname{Com}\left( \bbU_{m} \right) = \sup\{k \geq 1| \, f(k) \leq m\}.
    \end{align}
    Here, $\operatorname{Com}\left( \cdot \right)$ denotes the highest power exponent in the computational complexity, indicating that $\bbU_{m}^{\mathsf{all}}$ has complexity $O\left(n^{\operatorname{Com}\left( \bbU_{m}^{\mathsf{all}} \right)}\right)$.
\end{theorem}

\begin{proof}
    The proof consists of two main parts:
    \begin{enumerate}
        \item First, we demonstrate that the term of highest computational complexity emerges when order $m \ge f(k)$ (see Theorem~\ref{thm:bound_2_largest_numbers}).
        \item Second, we establish that all other terms exhibit complexity less than or equal to $n^k$ (see Theorem~\ref{thm:bound_2_contraction}).
    \end{enumerate}
\end{proof}


\begin{theorem}[Enumeration of Terms with Maximum Computational Complexity]
\label{thm:bound_2_largest_numbers}
\end{theorem}
\begin{proof}
   This problem can be reformulated as finding Eulerian paths in complete graphs. Specifically, we seek the minimum number of duplicate edges that must be added to a complete graph of order $k$ to obtain a graph containing an Eulerian path.

 By Euler's Theorem \citep{polya1983hamiltonian} for undirected graphs, a connected graph contains an Eulerian path if and only if either all vertices have even degree or exactly two vertices have odd degrees.
   
   In a complete graph of order $k$, each vertex has degree $k-1$ as it connects to all other vertices. We consider two cases:
   
   When $k$ is odd, the degree of each vertex ($k-1$) is even, thus the graph inherently contains an Eulerian path without requiring additional edges.
   
When $k$ is even, each vertex has odd degree ($k-1$). To construct an Eulerian path, we must add at least $\frac{k}{2}-1$ edges to ensure exactly two vertices maintain odd degree while all others attain even degree. This is achievable by adding the edges $\left(v_1, v_2\right), \left(v_3, v_4\right), \cdots, \left(v_{k-3}, v_{k-2}\right)$, which results in only vertices $v_{k-1}$ and $v_k$ having odd degree.

   This analysis leads to our function $f(k)$:
   \begin{align}
       f(k) = 
       \begin{cases}
           \binom{k}{2} + \frac{k}{2} - 1, & \text{if } k \text{ is even}, \\
           \binom{k}{2},                    & \text{if } k \text{ is odd},
       \end{cases}
   \end{align}
   where $\binom{k}{2}$ represents the initial number of edges in the complete graph, and the additional term $\frac{k}{2} - 1$ accounts for the necessary duplicate edges in the even case.
\end{proof}
\begin{theorem}[complexity upper bound]
\label{thm:bound_2_contraction}

\end{theorem}

\putbib[Master]
\end{bibunit}

\end{document}